\documentclass[11pt,twoside]{book}
\usepackage[T1]{fontenc}
\usepackage[utf8]{inputenc}
\usepackage{microtype}
\usepackage{geometry}
\geometry{
    paperwidth=6in,
    paperheight=9in,
    margin=1.1in,
    headheight=14pt
}
\usepackage{lettrine}
\usepackage{titlesec}
\usepackage{fancyhdr}
\usepackage{enumitem}
\usepackage{tabularx}
\usepackage{booktabs}
\usepackage{longtable}
\usepackage{tcolorbox}
\usepackage{adforn}
\usepackage{hyperref}
\usepackage{setspace}
\usepackage{parskip}

% Colors
\usepackage[dvipsnames]{xcolor}
\definecolor{royalblue}{RGB}{65,90,140}
\definecolor{agedpaper}{RGB}{250,245,235}

% Typography - A classic Garamond for that legal/manuscript feel
\usepackage{ebgaramond}
\renewcommand{\rmdefault}{ebgaramond}

% Chapter styling - formal, with rules
\titleformat{\chapter}[display]
{\normalfont\centering\Huge\bfseries\color{royalblue}}
{\chaptertitlename\ \thechapter}{20pt}{\Huge}
\titlespacing*{\chapter}{0pt}{0pt}{40pt}

% Section styling
\titleformat{\section}
{\normalfont\Large\bfseries\scshape}{\thesection}{1em}{}
\titleformat{\subsection}
{\normalfont\large\itshape}{\thesubsection}{1em}{}

% Header/footer
\pagestyle{fancy}
\fancyhf{}
\fancyhead[LE]{\thepage\quad\textsc{\nouppercase{\leftmark}}}
\fancyhead[RO]{\textsc{\nouppercase{\rightmark}}\quad\thepage}
\fancyhead[RE,LO]{}
\renewcommand{\headrulewidth}{0.4pt}
\renewcommand{\footrulewidth}{0pt}

% Custom environments
\newenvironment{sidebar}[1][]
{\begin{tcolorbox}[colback=agedpaper,colframe=royalblue,boxrule=0.5pt,arc=2mm,title={\scshape\small #1},fonttitle=\bfseries]}
{\end{tcolorbox}}

\newenvironment{quotebox}
{\begin{quote}\small\itshape}
{\end{quote}}

% List styling
\setlist[itemize]{noitemsep,topsep=0.5em}
\setlist[enumerate]{noitemsep,topsep=0.5em}

% Hanging indents for legal-style lists
\newcommand{\legalitem}[1]{\par\noindent\textbf{#1}\par\noindent\hangindent=2em\hangafter=1}

\title{\Huge\scshape The Hedge-Law Realm\\\vspace{0.5em}\Large\normalfont\textit{A Viterran Expansion for Fate's Edge}}
\author{\textit{Set down by Dusana of the Raven Road, under a borrowed lamp in the White Hall}\\[0.5em]\small\textit{The ink is gall and the ash of a burned writ. The binding is leather from a debtor's cow.}}
\date{}

\begin{document}

\frontmatter

\maketitle

\thispagestyle{empty}
\begin{center}
\vspace*{2em}
\textit{The ink is gall and the ash of a burned writ.}\\
\textit{The binding is leather from a debtor's cow.}
\vspace{2em}
\end{center}

\tableofcontents

\mainmatter

\chapter*{Preface: The Queen's Peace}
\addcontentsline{toc}{chapter}{Preface: The Queen's Peace}

\lettrine[lines=3,findent=2pt]{T}{he law is not a sword.} It is a scale. A sword can be stolen; a scale only tilts.

\begin{quotebox}
--- Queen Veralyn Everblood, from her first address to the Parlement
\end{quotebox}

Viterra is not a land of crumbling ruins or haunted moors. It is a land of hedgerows and courthouses, of ancient customs written in margins, of disputes settled by gavel and precedent rather than by blade. The Queen's Peace is not a promise---it is a contract, renewed each morning by every magistrate, every justicar, every clerk who opens a ledger.

After the succession wars that bled the eastern marches for a generation, Queen Veralyn Everblood sits on a throne still warm from her rivals' blood. She is young, unmarried, and surrounded by counselors who mistake caution for loyalty. Her reign is defined by three pressures:

\begin{itemize}
    \item The need to rebuild law where war left only custom.
    \item The presence of turncloak nobles who switched sides multiple times during the wars and now demand pardon.
    \item The quiet power of the Justicars---magistrates who can travel anywhere, hear any case, and speak with the Queen's own voice.
\end{itemize}

This expansion provides the tools to play in Viterra: legal duels, oath-courts, hedge-law, and the politics of a realm that has learned, at great cost, that justice is the only thing that holds a kingdom together when armies have gone home.

\chapter{The Pillars of Viterran Law}

Viterran law is not written in a single codex. It lives in four sources, each with its own weight.

\section{The Queen's Word}

The monarch may issue decrees, but a decree that contradicts the \textit{Codex of Precedents} is null unless confirmed by the Parlement. Queen Veralyn has issued only seven decrees in five years; her predecessor issued forty in a single session.

A royal decree carries the weight of the crown's authority. To defy it is treason. But the Queen herself may be petitioned to set aside a decree that causes manifest harm---a check on absolute power that has not been used in living memory.

\section{The Codex of Precedents}

A massive book, kept in the White Hall under lock and witness, recording every judgment of the Queen's Justicars for four centuries. When a new case arises, the first question is always: \textit{"What does the Codex say?"} If no precedent exists, the Justicar creates one---and the next Justicar must follow it.

The Codex is not infallible. Pages have been lost to fire, water, and deliberate excision. Some precedents contradict one another; older clerks whisper that the Codex sometimes \textit{changes} overnight, a new clause appearing in a margin as if written by a ghost.

\begin{sidebar}[Strings (Viterran Law)]
\begin{itemize}
    \item \textbf{Precedent Citation} -- Reroll a failed Insight roll related to law or custom.
    \item \textbf{Codex Page} -- A copied folio; once per session, reduce DV by 1 on a legal argument.
\end{itemize}
\end{sidebar}

\section{The Hedge-Law}

The unwritten customs of villages and rural manors. A hedge-court has no gavel, only a stool and a witness. Its judgments are binding only on those who consent---but refusal to consent is itself a mark against one's reputation. Hedge-law covers boundary stones, water rights, grazing rotations, and the proper distance between beehives.

Hedge-courts are held under ancient oaks or at crossroads. The judge is often the oldest person present, or the one with the most even temper. Verdicts are pragmatic: a stolen pig is replaced; a broken fence is mended by the one who broke it; a blood-feud is settled by a marriage.

\begin{sidebar}[Strings (Hedge-Law)]
\begin{itemize}
    \item \textbf{Boundary Stone} -- Declare a disputed line once; the hedge-court will honor your claim for a season.
    \item \textbf{Hedge Writ} -- Safe passage through rural lands; no tolls, no questions.
\end{itemize}
\end{sidebar}

\section{The Right of Acknowledgment}

The newest pillar, added by Queen Veralyn's own hand after the trial of Baron Aldric the Turncloak. It permits any child who was not ransomed, not visited, or not named in a will to demand that a parent speak one truth aloud under witness.

The truth need not be kind. It need only be honest.

Once spoken, the debt is considered paid in full. The parent cannot be compelled again. The child may walk away---or stay. The law does not care.

Mechanically, invoking the Right of Acknowledgment is a social action (Presence + Sway or Command) against the parent's Resolve. On a hit, the parent speaks a true answer to one question. On a critical, the answer must be complete and without evasion. The parent may refuse, but refusal is an admission of guilt in the court of public reputation---they become \textit{anima non grata}, and no lawful burial ground will accept their body.

\chapter{Factions of Viterra}

\section{The Queen's Justicars}

A small order of traveling magistrates, each bearing a silver gavel and a sealed writ. A Justicar may hear any case, anywhere, and her judgment is final unless overturned by the Queen herself. There are only twelve Justicars in all of Viterra. They are feared, respected, and occasionally murdered.

\begin{sidebar}[Strings]
\begin{itemize}
    \item \textbf{Gavel token} (force a parley under witness)
    \item \textbf{Justicar's seal} (DV -1 for legal negotiations)
    \item \textbf{Precedent citation} (reroll a failed Insight roll related to law)
\end{itemize}
\end{sidebar}

\section{The Hedge-Courts}

Local assemblies of elders, millers, and retired soldiers who settle disputes without reference to the Codex. Their judgments are pragmatic, swift, and often brutal. A hedge-court can outlaw a person from a village---a sentence of slow starvation.

\begin{sidebar}[Strings]
\begin{itemize}
    \item \textbf{Boundary stone} (declare a disputed line once)
    \item \textbf{Hedge writ} (safe passage through rural lands)
    \item \textbf{Goose feather} (one use to call a moot)
\end{itemize}
\end{sidebar}

\section{The Turncloak Nobility}

Nobles who changed sides during the succession wars. Some did so out of opportunity; others out of fear. Queen Veralyn has pardoned most but forgotten none. Their estates are intact, but their honor is stained. Many seek redemption through service---or revenge through conspiracy.

\begin{sidebar}[Strings]
\begin{itemize}
    \item \textbf{Pardon letter} (reduces Exposure by 2 in Viterra)
    \item \textbf{Oath of fealty} (once/arc, call on a turncloak for aid, but mark a Debt to them)
    \item \textbf{Blood-silver receipt} (proof of payment for a past betrayal)
\end{itemize}
\end{sidebar}

\section{The Parlement of Viterra}

An assembly of nobles, bishops, and guildmasters that meets twice a year to advise the Queen and approve taxes. The Parlement has no power to make law, but it can refuse to fund the Queen's projects. Veralyn has learned to charm, bribe, and threaten them by turns.

\begin{sidebar}[Strings]
\begin{itemize}
    \item \textbf{Proxy vote} (once/session, sway a Parlement decision)
    \item \textbf{Tax receipt} (reduces Supply cost by 1 in Viterran markets)
    \item \textbf{Censure motion} (public shaming of an enemy)
\end{itemize}
\end{sidebar}

\chapter{New Talents: The Justicar's Path}

For characters who walk the road of law, oath, and judgment. These talents are available to any class, but they resonate most strongly with Runekeepers of Mykkiel or The Witness.

\section{Foundation Talents (2--4 XP)}

\legalitem{Cite Precedent (2 XP)} Once per session, you may declare that a past judgment covers the current situation. Roll Wits + Lore (DV 3). On a hit, gain +1 die on your next legal or social roll. On a miss, the precedent is contested---GM gains 1 SB.

\legalitem{Gavel's Authority (4 XP)} You carry a small gavel (or a symbolic equivalent). When you strike it against a hard surface and state a judgment, hostile NPCs within Near range must test Resolve (DV 4) or hesitate for one exchange. Usable once per scene.

\section{Advanced Talents (6--8 XP)}

\legalitem{Right of Acknowledgment (6 XP)} \textit{Requirement: You have witnessed an injustice involving kin or oath.} Once per arc, you may demand that an NPC speak one true answer to a question of your choice. The NPC may refuse, but refusal forfeits all legal protections (they become \textit{anima non grata} in Viterran territory). The truth need not be complete, but it must not be a lie.

\legalitem{Justicar's Immunity (8 XP)} You have been granted the Queen's own seal. You cannot be arrested or detained by any lesser authority in Viterra. However, abuse of this immunity (committing a crime in plain sight) will cause the seal to crack---lose this talent until you undergo a ritual of purification (a downtime project).

\chapter{Patrons of the Viterran Courts}

\section{The Witness (Viterran Variant)}

Here, The Witness does not reveal hidden truths---it \textit{attests} to sworn testimony. Its shrines are courthouses; its rites are oaths spoken under a single lamp.

\legalitem{Low Rite --- Testimony (4 XP)} When you repeat a statement you heard under oath, you may roll Presence + Investigation (DV 3). On a hit, you know whether the statement was false. (Mark +1 Obligation)

\legalitem{Standard Rite --- Unbreakable Vow (8 XP)} You and a willing partner swear an oath before witnesses. Until the oath is fulfilled, both gain +1 die to any action that serves it. Breaking the oath causes 2 Harm (Spirit). (Mark +2 Obligation)

\section{Mykkiel (Judge of the Scar)}

Mykkiel's Viterran face is that of a circuit judge, riding from town to town, carrying a ledger of unpaid debts. His followers are Justicars who have lost their faith in mercy but not in order.

\legalitem{Low Rite --- Summons (4 XP)} Name a person within Far range. They hear your voice as if you stood beside them. They are not compelled to obey, but they know the message is legal. (Mark +1 Fatigue)

\legalitem{Standard Rite --- Attachment (8 XP)} Touch an object or person. For one day, any attempt to move them from their current location suffers +1 DV. (Mark +2 Obligation)

\section{The Queen (Mortal Patron)}

Queen Veralyn herself can serve as a Patron for those sworn directly to her service. Her Rites are not magical---they are legal and social permissions.

\legalitem{Low Rite --- Royal Favor (4 XP)} Present a token with the Queen's seal. For one scene, any Viterran official treats you as if you were a royal messenger. (Mark +1 Exposure -- the Queen's enemies take note)

\legalitem{Standard Rite --- Pardon (8 XP)} Forgive one crime committed by another person, as long as that crime did not harm the Queen's interests. The pardon is legally binding. (Mark +2 Obligation to the Queen; she may call on you later)

\chapter{The Justicar's Judgment}

\begin{quotebox}
\textit{The tale of Baron Aldric the Turncloak, set down by Dusana of the Raven Road after a cup of small beer in a Viterran hedge-court. The ink is gall and the memory of a father's silence.}
\end{quotebox}

\section{The Charges}

In the fifth year of Queen Veralyn's peace, the Queen's Justicar---a woman named Elara Venn---received a sealed writ. The wax was black, stamped with the crown of Viterra and the broken column of Vhasia. Inside, a single sheet of vellum:

\begin{quote}
\textit{You are commanded to hear the case of Baron Aldric of the Scar, formerly Sir Aldric of the White Moor, accused of treason thrice sworn and thrice forsworn, of oath-breaking before witnesses, and of refusing the lawful ransom of his own blood. Let the ledger be balanced.}
\end{quote}

Elara read the name. Her hand did not shake.

She had squired for Aldric twenty years ago, when he was a young knight with a quick laugh and a quicker sword. She had carried his shield at the Ford of Ashes. She had watched him bargain for his first holding---a burnt tower and a handful of serfs---with the same smile he used to haggle over a horse.

She had also watched him abandon Queen Veralyn's mother at the Battle of the Three Bridges, switching his banner to the Pretender's side an hour before the charge. When the Pretender lost, Aldric swore fealty again. Elara had testified on his behalf at the first inquiry.

She was young then. She believed in second chances.

\section{The Knight's Ledger}

The court convened in the White Hall of Viterra, a long room lined with oak and the painted faces of dead magistrates. Baron Aldric was brought in chains of gilded bronze---a courtesy, because he still had friends. His hair had gone grey. His eyes had not.

Elara read the first count: \textit{That you swore loyalty to Queen Veralyn's father, then to the usurper Caster, then again to the Queen's father, then to the Duke of Adria, and finally to the Queen herself, each time breaking the previous oath without payment of blood-silver.}

Aldric smiled. "A man must survive, Justicar. The throne changed hands. I changed my hand. Is that not the custom of Viterra?"

The hall murmured. Elara did not smile.

"The custom of Viterra," she said, "is that an oath is a bridge. You do not burn a bridge behind you unless you have paid the toll in full. You owe eighteen years of back-tolls, Aldric. That is not survival. That is usury."

She let the silence stretch. Then she read the second count---the one that had kept her awake.

\section{The Son Who Was Not Ransomed}

"In the year of the Broken Crown, after the Battle of the Red Ditch, your son---Ser Roran, then seventeen---was captured by the forces of the Duke of Adria. The customary ransom for a knight's eldest son was three hundred silver marks. The Duke's herald offered you a discount of fifty marks if you paid within the week."

Aldric's smile thinned. "I recall."

"You refused. You said, and I quote the herald's sworn testimony: `I have more sons and can make further. The boy was weak. Let the Duke keep him.'"

A rustle passed through the hall. Elara did not look at the gallery, where a lean, grey-faced man sat in the shadows---Ser Roran, now thirty-seven, his left hand sheathed in a gauntlet of silvered steel, his right hand resting on the pommel of a sword that had seen twenty campaigns.

"The Duke did not keep him," Elara continued. "He threw the boy into a debtor's cell, expecting you to relent. You did not. For three years, Roran rotted. Then, a stranger paid the ransom---a hedge-knight who had once served under your banner and remembered the boy's courage. The knight gave every coin he had. He died of a fever the following winter. Roran took his name and his sword."

She paused.

"He did not return to you. He fought in the eastern wars. He won the Queen's own pennon at the Siege of the Horn. He knighted himself with a broken blade and the blessing of a battlefield priest. He has not spoken your name in seventeen years."

Aldric's jaw tightened. "The boy survived. Stronger for it."

Ser Roran stood. The gallery fell silent. He walked to the bar and placed his gauntleted hand on the rail.

"Father," he said, "do you know what I remember most? Not the cell. Not the chains. The moment I heard you say it. I was hiding behind the herald's tent. I heard you laugh."

\section{The Justicar's Judgment}

The law of Viterra is ancient and precise. Treason is punishable by death. Oath-breaking forfeits all lands. Refusing a son's ransom is not a crime---it is merely cruelty.

Elara Venn sat on the bench for three hours. She read the old charters. She consulted the \textit{Codex of Precedents}, a book so heavy that two clerks carried it. She asked for a cup of water. She asked for silence.

When she spoke, her voice was level.

"Baron Aldric of the Scar, on the first count---treason by repeated oath-breaking---I find you guilty. The sentence is forfeiture of all lands, titles, and pensions. You will be stripped of your banner and your ring. You will live as a hedge-knight if any will take you."

Aldric's face did not change. He had expected this.

"On the second count," Elara continued, "refusal of ransom and abandonment of kin---the law is silent. Viterra does not compel a father to love. But the law does have a clause for debts of the heart. It is called the \textit{Right of Acknowledgment}. No one has invoked it in two hundred years."

She looked at Ser Roran.

"Ser Roran, son of Aldric, you have the right to demand that your father speak one truth aloud, in this court, under witness. If he refuses, he is declared \textit{anima non grata}---his soul outlawed from all Viterran burial grounds. His body will be buried at a crossroads with a stake through his chest. The law does not forbid it. The Church does not approve. But the law is the law."

Aldric laughed. "A truth? I have told a thousand truths. I have told the truth that a man must cut what he cannot carry. I have told the truth that sons are arrows---you loose them and you do not weep. That is my truth."

"No," Ser Roran said. His voice was soft. "Your truth is that you were afraid. Not of the Duke. Of me. I was stronger than you at seventeen. You saw it. You could not bear to pay my ransom because then you would owe me, and you have never owed anyone in your life."

Aldric opened his mouth. Closed it.

The hall waited.

"I am not asking you to love me," Roran said. "I am asking you to say it. Just once. Say: `I was afraid.'"

Aldric stared at his son. For a long moment, the mask of the turncloak---the smile, the shrug, the armor of indifference---cracked.

"I was afraid," he whispered.

The sound was so small that the clerks leaned forward.

"Of what?" Elara asked.

"Of him. Of what he would become. He would have been better than me. And I\ldots{} I could not stand to watch."

Ser Roran turned away. He did not weep. He had no tears left. He walked to the door, his silvered gauntlet clicking against the stone.

Elara raised her gavel.

"The Right of Acknowledgment is satisfied. The court is adjourned. Baron Aldric---former Baron Aldric---you are free to go. But I advise you to leave Viterra by sunset. The road is long, and the winter is coming."

\section{The Aftermath}

Aldric died three years later in a ditch outside Thepyrgos. No one claimed the body. The crossroads stake was not needed; he had outlived his own burial rights.

Ser Roran became a knight of the Queen's household. He wore the silver gauntlet over his left hand---not to hide the damage, but to remind himself that a stranger's coin had bought his life. He never took ransom. Once, a captured lord offered him five hundred marks for his freedom. Roran refused.

"Pay it to the hedge-knight's widow," he said. "I have no need for gold. I had a father once. That was expensive enough."

Queen Veralyn, when she heard the verdict, said nothing for a long time. Then she called for the \textit{Codex of Precedents} and wrote a new clause in the margin, in her own hand:

\begin{quote}
\textit{"The Right of Acknowledgment may be invoked by any child who was not ransomed, not visited, not named in a will. The truth, once spoken, is payment in full. No further interest shall accrue."}
\end{quote}

The clerks copied the clause into every new edition. It is still law.

\begin{center}
\textit{The ink is dry. The gavel is still.}
\end{center}

\chapter{Adventure Seeds}

\begin{enumerate}
    \item \textbf{The Unfinished Testament.} A dying noble calls for a Justicar to witness his will. The will disinherits his eldest son in favor of a bastard daughter who was never ransomed from captivity. The son hires thugs to intercept the Justicar. The party must deliver her to the bedside before the old man dies.
    
    \item \textbf{Hedge-Court at Dusk.} A village accuses a turncloak knight of stealing a boundary stone. The knight claims the stone was always on his land. The party must survey the hedge, interview the oldest residents, and render a verdict that will not start a blood-feud.
    
    \item \textbf{The Turncloak's Daughter.} A Vhasian baron's daughter, now living in Viterra, petitions for the Right of Acknowledgment against her father's estate. The father is dead, but his executors refuse to release her inheritance. The party must find a witness who heard the father speak a damning truth---or forge a confession that the hedge-court will accept.
    
    \item \textbf{The Codex Theft.} A page from the \textit{Codex of Precedents} has been stolen---a page that would decide a pending lawsuit between two powerful houses. The party must track the thief through the underworld of legal scribes and forgers before the case is dismissed for lack of precedent.
    
    \item \textbf{The Queen's Suitors.} Queen Veralyn has announced she will consider marriage offers. Every noble house in Viterra sends a champion to court her. The party is hired by a minor house to sabotage the other suitors---legally, of course. No poison. No daggers. Just clever readings of dowry law, ambush by precedent, and a well-timed Right of Acknowledgment.
\end{enumerate}

\chapter{GM Toolkit: Clocks and Fronts}

\section{Sample Clocks}

\begin{longtable}{@{}p{2.5cm}p{4cm}p{5cm}@{}}
\toprule
\textbf{Clock} & \textbf{Purpose} & \textbf{Advance / Retreat} \\
\midrule
\endhead
\textbf{The Queen's Patience [6]} & How long the crown tolerates noble infighting & Advance when parties flout Justicar rulings; retreat when a major case is settled cleanly \\
\addlinespace
\textbf{Hedge-Law Fracture [4]} & Local customs breaking down under pressure & Advance when a hedge-court's verdict is ignored; retreat when a respected elder restores peace \\
\addlinespace
\textbf{Turncloak Conspiracy [8]} & A plot among pardoned nobles to reclaim lost lands & Advance when the party ignores rumors or protects a turncloak; retreat when a conspirator is exposed \\
\addlinespace
\textbf{Codex Decay [6]} & The \textit{Codex} losing authority as precedents contradict & Advance when contradictory rulings are issued; retreat when a Justicar reconciles two precedents \\
\bottomrule
\end{longtable}

\section{Sample Front: The Succession Crisis}

\textbf{Impulse:} The Queen has no heir. Her marriage negotiations have stalled. The turncloak nobility smell blood.

\textbf{Grim Portents:}
\begin{itemize}
    \item A hedge-court declares a local lord has the "ancient right" to name the next Queen.
    \item The Parlement refuses to fund the Justicars, leaving a key case unheard.
    \item A turncloak knight invokes the Right of Acknowledgment against the Queen herself, demanding she name her father's true heir.
\end{itemize}

\textbf{Clocks:}
\begin{itemize}
    \item \textit{Throne Vacancy [8]} -- When full, a claimant marches on the White Hall.
    \item \textit{Justicar Defection [4]} -- When full, one of the twelve Justicars declares for a rival.
    \item \textit{People's Petition [6]} -- When full, a hedge-court of commoners issues a binding judgment on the succession, overriding noble law.
\end{itemize}

\chapter{Colophon}

This expansion was written in a Viterran hedge-court, on paper that had been used to wrap cheese. The ink is gall and the ash of a burned writ. The binding is leather from a debtor's cow.

The Queen's Justicar Elara Venn died in her sleep seven years after the trial of Baron Aldric. Her gavel is kept in the White Hall, next to the \textit{Codex of Precedents}. Children who visit the Hall are told to touch it for luck. They do.

Ser Roran served the Queen for thirty years. He never married. He never took a squire. He wore the silver gauntlet until the day he died, and it was buried with him. The hedge-knight who ransomed him has no grave---his body was never found. But Roran paid his widow every year on the anniversary of the ransom, in full, with interest.

The road to Viterra is straight, the hedges are trimmed, and the law is patient.

\begin{center}
\vspace{2em}
\textit{The ink is dry. The gavel is still.}
\end{center}

\end{document}
